\scanpage{113}
-- et quand les composés (158) sont connus, on vérifie que les morphismes $D_J-D_{J'}$ sont connus,\ldots

On est donc intéressé plus particulièrement par les hom
\[
(164)\qquad D_{d,\delta}\unsure{-|\longrightarrow}D_\delta\qquad(d>\delta)
\]
qui jouent le rôle des ``morphismes but'' alors que les $D_{d,\delta}-D_d$ jouent le rôle des ``morphismes source'', et $D_{d,\delta}$ le rôle de ``\deleted{\unsure{f}}multiplicité des flèches'' (\unsure{ou|où} mor. d'inclusion) -- entre éléments de sous-multiplicités \inserted{de $X$} associés à $D$.

On a des isom. can.
\[
(165)\qquad D_{d,d'}\times_{D_{d'}}D_{d',d''}\simeq D_{d,d',d''}\qquad\text{si }d>d'>d''
\]
et on \deleted{\unsure{en}} déduit de cette description les morphismes canoniques
\diagram{
\begin{tikzcd}[ampersand replacement=\&, column sep=large,row sep=small]
\& \& D_d\\
\& D_{d,d'}\arrow[ur]\arrow[dr] \&\\
D_{d,d',d''}\arrow[ur]\arrow[dr,dash] \& \& D_{d'}\\
\& D_{d',d''}\arrow[ur]\arrow[dr] \&\\
\& \& D_{d''}
\end{tikzcd}
\qquad(166)
}
mais non directement \deleted{$\unsure{D_{d'}}$}
\[
(167)\qquad D_{d,d',d''}\longrightarrow D_{d,d''},
\]
mais on note que
\[
(168)\qquad D_{d,d''}\unsure{-|\longrightarrow}D_d\times D_{d''}.
\]
\margin{112}
